\section{When does aggregation pay: analysis}\label{sec:pay}

Fewer iterations are not the same as less work. This section and
Section~\ref{sec:pay-meas} separate the two. Here: an exact expression
for how much more progress a block makes than the single MDM step at
the same iterate, a cost model for what the block costs, and a theorem
that bounds the drop steps by an identification phase rather than by
the crude count of Theorem~\ref{thm:linear}. Section~\ref{sec:pay-meas}
measures every quantity in the model on the five real cells of
Section~\ref{sec:exp}.

\subsection{Progress per unit of work at one iterate}

Consider a case-(a) iteration with $k'$ retained pairs, all uncapped,
so that $\gamma_j=\mathrm{gap}_j/\norm{d_j}^2$ and pair $j$ alone would
gain $\Delta_j=\mathrm{gap}_j^2/(2\norm{d_j}^2)$, with $\Delta_1$ the MDM
gain. Write
\[
Q=\frac{\sum_{j\le k'}\Delta_j}{\Delta_1},\qquad
\chi=\frac{D}{S},\qquad
S=\sum_j\gamma_j\mathrm{gap}_j=2\sum_j\Delta_j,\qquad
D=\Big\lVert\sum_j\gamma_jd_j\Big\rVert^2 .
\]
$Q$ measures how much useful progress the block has beyond its leading
pair and $\chi$ how much the directions lose (or gain) by being
combined; orthogonal directions give $\chi=1$. Call the block
\emph{$\eta$-diverse} if $|\ip{d_i}{d_j}|\le\eta\norm{d_i}\norm{d_j}$ for
$i\ne j$ and \emph{$\beta_{\mathrm{flat}}$-flat}, $0<\beta_{\mathrm{flat}}\le1$, if
$\Delta_j\ge\beta_{\mathrm{flat}}\Delta_1$ for every $j$ (the subscript keeps
the flatness parameter apart from the $W$-side weights $\beta$).

\begin{proposition}[Exact gain ratio]\label{prop:ratio}
At such an iteration the exact line search gives
\[
\frac{\Delta_B}{\Delta_1}=Q\cdot\begin{cases}2-\chi,&\chi\le1,\\
1/\chi,&\chi\ge1.\end{cases}
\]
If the block is $\eta$-diverse then $\chi\le1+\eta(k'-1)$; if it is
$\beta_{\mathrm{flat}}$-flat then $Q\ge1+\beta_{\mathrm{flat}}(k'-1)$; hence the guarded step gains at
least
\[
\max(\Delta_B,\Delta_1)\ \ge\ \Delta_1\,\frac{1+\beta_{\mathrm{flat}}(k'-1)}{1+\eta(k'-1)} .
\]
\end{proposition}
\begin{proof}
For uncapped pairs $\gamma_j\norm{d_j}^2=\mathrm{gap}_j$, so
$S=\sum_j\mathrm{gap}_j^2/\norm{d_j}^2=2\sum_j\Delta_j$ and
$\varphi_j^2=\gamma_j^2\norm{d_j}^2=\gamma_j\mathrm{gap}_j$, whence
$\sum_j\varphi_j^2=S$. The line search takes $\theta=\min(S/D,1)$. If $D\ge S$
the optimum is interior and $\Delta_B=S^2/(2D)=(S/2)/\chi=\sum_j\Delta_j/\chi$;
if $D<S$ then $\theta=1$ and $\Delta_B=S-D/2=(S/2)(2-\chi)$. Dividing
by $\Delta_1$ gives the two cases. For the bound on $\chi$,
$D=\sum_{i,j}\gamma_i\gamma_j\ip{d_i}{d_j}\le\sum_j\varphi_j^2+\eta\sum_{i\ne j}\varphi_i\varphi_j
\le(1+\eta(k'-1))\sum_j\varphi_j^2=(1+\eta(k'-1))S$ by Cauchy--Schwarz on the
cross terms; flatness gives $\sum_j\Delta_j\ge\Delta_1(1+\beta_{\mathrm{flat}}(k'-1))$.
When $\chi\le1$ the ratio is $Q(2-\chi)\ge Q$, and when $\chi\ge1$
it is $Q/\chi$; both are at least the stated bound.
\end{proof}
Proposition~\ref{prop:ratio} recovers the diversity bound of Proposition~\ref{prop:orth} and adds the comparison with $\Delta_1$. The two
numbers $Q$ and $\chi$ are computed by the guard at every block step
(from $S$, $D$ and the per-pair gaps), so both are observable at no
extra cost. Note the direction of the inequalities: the theorem's $\eta$
bounds the \emph{maximum} absolute pairwise cosine of the retained
directions, whereas the \emph{effective correlation}
\[
\hat\eta:=\frac{\chi-1}{k'-1}\qquad(k'>1)
\]
that a block step reveals is a $\gamma$-weighted aggregate in which
cross terms can cancel; the proof gives $\hat\eta\le\eta$, and $\hat\eta$
can be negative. Likewise $Q/k'$ close to $1$ says the average pair is
as good as the leading one, not that every pair is ($\beta_{\mathrm{flat}}$-flatness).
The empirical model of Section~\ref{sec:pay-meas} is therefore stated in
$Q$ and $\chi$, the quantities the guard measures, and the theorem's
assumptions (all retained pairs uncapped, maximum cosine at most $\eta$,
minimum relative gain at least $\beta_{\mathrm{flat}}$) are recorded separately.

\paragraph{Cost.} Per iteration the solver pays a scan and certificate
cost that does not depend on $k$, a step cost that grows with the number
of retained pairs (the $k'$-column score update and the $O(k'^2)$ guard), and one Gram or kernel column for every
receiver or donor not in the cache. Write
\[
\mathcal C_B=a+b\,k'+c_{\mathrm{col}}\,n^{\mathrm{new}}_B,\qquad \mathcal C_1=a+b+c_{\mathrm{col}}\,n^{\mathrm{new}}_1,
\]
for the block and the single step from the same state, where $n^{\mathrm{new}}_B\ge n^{\mathrm{new}}_1$
count the column misses of the two candidates ($n^{\mathrm{new}}_1\le1$: only the MDM
receiver can be new) and $a,b,c_{\mathrm{col}}$ are timing coefficients.
They are fitted, not assumed: they depend on $n$, $d$, the
representation (dense, kernel) and the hardware, and
Section~\ref{sec:pay-meas} fits them per cell. The model is a timing
approximation of the implementation, not an upper bound on
computational work. The guarded block improves progress per unit of work at this
iterate exactly when
$\max(\Delta_B,\Delta_1)/\Delta_1>\mathcal C_B/\mathcal C_1$; by
Proposition~\ref{prop:ratio} a sufficient condition is
\begin{equation}\label{eq:ppc}
\frac{1+\beta_{\mathrm{flat}}(k'-1)}{1+\eta(k'-1)}\;>\;\frac{a+bk'+c_{\mathrm{col}}n^{\mathrm{new}}_B}{a+b+c_{\mathrm{col}}n^{\mathrm{new}}_1}.
\end{equation}
Two regimes are visible in \eqref{eq:ppc}. If $c_{\mathrm{col}}$
dominates $a$ and $b$, the right-hand side tends to the miss ratio
$n^{\mathrm{new}}_B/n^{\mathrm{new}}_1\le k'$: the block must earn every new column it
touches beyond the single step's, and wins by the full gain ratio when
both need the same column ($n^{\mathrm{new}}_B=n^{\mathrm{new}}_1$), so column dominance
raises the hurdle on the iterations on which the block's extra receivers
are not cached; it does not exclude a gain. If the per-iteration terms dominate, the
right-hand side is $(a+bk')/(a+b)$, linear in $k'$, while the left-hand
side saturates at $\beta_{\mathrm{flat}}/\eta$; the ratio of the two is maximised at
\begin{equation}\label{eq:kstar}
k^*=\sqrt{\frac{a(1-\eta)}{b\,\eta}}\qquad(\beta_{\mathrm{flat}}=1),
\end{equation}
the block size at which the marginal cost of one more pair equals its
marginal gain. Equation~\eqref{eq:kstar} is the unconstrained continuous
optimiser of the model for $a,b>0$ and $0<\eta<1$: as $\eta\to0$ it grows
without bound (every pair helps), as $\eta\to1$ it tends to $0$ (only the
leading pair helps), and an implementation rounds it and clips it to
$\{1,\dots,k_{\max}\}$.

\subsection{Drops are an identification cost}

Theorem~\ref{thm:linear} charges up to $k$ drop steps per non-drop
iteration, which would cancel a $k$-fold gain in progress. The
measurements below show that drops in fact happen while the support is
being built and stop once it has settled. The following theorem proves
this under the standard non-degeneracy assumptions of active-set
identification \citep{bomze2020,garber2020}: after an explicit
iteration, every step that removes a non-face point is a drop, no
non-face point enters, and once the supports lie in the optimal faces
and the weights are close to the optimal ones, no step is a drop at all.

Write $\mathcal F_V=\conv\{v_i:i\in F_V\}$ and
$\mathcal F_W=\conv\{w_j:j\in F_W\}$ for the geometric faces ($F_V,F_W$
remain index sets), and let $V_{F_V}\in\R^{d\times|F_V|}$,
$W_{F_W}\in\R^{d\times|F_W|}$ be the matrices whose columns are the face
points. A witness pair $(p^*,q^*)$ with $p^*=\sum_i\alpha^*_iv_i$,
$q^*=\sum_j\beta^*_jw_j$ and $p^*-q^*=z^*$ carries optimal weights
$\alpha^*,\beta^*$.
\begin{assumption}\label{ass:nondeg}
(A1) \emph{Strict complementarity}: $F_V=\operatorname{supp}\alpha^*$ and
$F_W=\operatorname{supp}\beta^*$, and
$\alpha_{\min}:=\min\big(\min_{i\in F_V}\alpha^*_i,\ \min_{j\in F_W}\beta^*_j\big)>0$
is the smallest positive optimal weight over both classes. (A2)
\emph{Non-degeneracy}: the witness pair $(p^*,q^*)$ is unique, the
vertices of each optimal face are affinely independent (so
$\alpha^*,\beta^*$ are unique), and $\sigma>0$ denotes a constant with
$\norm{\alpha-\alpha^*}_1+\norm{\beta-\beta^*}_1\le\norm{(p-q)-z^*}/\sigma$
for all $p\in\mathcal F_V$, $q\in\mathcal F_W$ with face weights
$\alpha,\beta$ (such a constant exists under the two preceding
conditions and is computed below). Let $d_{\min}>0$ be the smallest distance from a face
vertex to the affine hull of the other vertices of its face, the minimum
taken over the faces with at least two points; if both faces are
singletons, $(p^*,q^*)$ is the unique pair of face points and the
no-drop conclusions below hold as soon as the supports lie in the faces.
\end{assumption}
Under (A1) and (A2), affine independence and uniqueness of the witness
pair make the linear map $(\alpha-\alpha^*,\beta-\beta^*)\mapsto(p-q)-z^*$ injective on the
subspace $U=\{\sum_i(\alpha_i-\alpha^*_i)=0=\sum_j(\beta_j-\beta^*_j)\}$
of $\R^{|F_V|+|F_W|}$; if $\sigma_2>0$ is the smallest singular value of
$[V_{F_V},-W_{F_W}]$ restricted to $U$, then
$\sigma=\sigma_2/\sqrt{|F_V|+|F_W|}$ works, the square root being the
conversion from the Euclidean to the $\ell_1$ norm of the weights. When
$U=\{0\}$ (both faces singletons) any $\sigma>0$ works and the weight
condition is void.

The distance to the optimum is controlled by the objective error
directly: $z^*$ minimises the $1$-strongly convex $F$ over the convex
$Z$, so $\ip{z^*}{z-z^*}\ge0$ for every $z\in Z$ and
\begin{equation}\label{eq:dist-h}
\tfrac12\norm{z-z^*}^2\ \le\ \ip{z^*}{z-z^*}+\tfrac12\norm{z-z^*}^2\ \le\ h(z),
\qquad\text{i.e.}\quad \norm{z_t-z^*}\le\sqrt{2h_t}.
\end{equation}
(The screening radius $\varrho_t\ge\norm{z_t-z^*}$ is the observable
proxy for this distance; the theorem does not need it.) Let $\ell_t$ be
the total support size, $\zeta_k=k+1$ (one-sided) or $2k+1$ (two-sided
schedule), and define
\begin{align*}
H_{\mathrm{id}}&=\min\Big\{\frac{\tau^2}{72L^2},\ \frac{\tau}{12},\ h_3\Big\},\qquad
h_3=\max\{h:\ 2h+R\sqrt{2h}\le\tau/6\},\\
H_{\mathrm w}&=\min\Big\{H_{\mathrm{id}},\ \frac{(\sigma\alpha_{\min})^2}{8},\
\frac{(\alpha_{\min}d_{\min})^2}{8}\Big\}.
\end{align*}

\begin{theorem}[Identification and drop count]\label{thm:ident}
Let $\delta^*>0$, let $\tau,L,R$ be as in Theorem~\ref{thm:activation}
and let the guarded block-transfer algorithm run with any $k\ge1$ under
either schedule of Theorem~\ref{thm:linear}. Let $t_{\mathrm{id}}$ be the
first iteration with $h_t<H_{\mathrm{id}}$ and $N$ the number of
non-face support indices at $t_{\mathrm{id}}$.
\begin{enumerate}\setlength\itemsep{0pt}
\item[(i)] For every $t\ge t_{\mathrm{id}}$: if either side's support
contains a non-face point, the step taken is a drop of a non-face point;
and no non-face point enters either support. Hence
$\operatorname{supp}\alpha\subset F_V$ and $\operatorname{supp}\beta\subset F_W$
from iteration $t_{\mathrm{id}}+N$ on, and
$N\le \ell_{t_{\mathrm{id}}}\le \ell_0+\zeta_k\,t_{\mathrm{id}}$.
\item[(ii)] Under Assumption~\ref{ass:nondeg}, let $t_{\mathrm w}$ be the
first iteration with $h_t<H_{\mathrm w}$. No iteration
$t\ge\max(t_{\mathrm w},t_{\mathrm{id}}+N)$ is a drop, and no block pair
is capped at such an iteration. Consequently every drop occurs before
$t_{\mathrm w}+N$, the number of drops is at most
$\zeta_k(t_{\mathrm w}+N)+\ell_0$, and
$t_{\mathrm w}\le \ell_0+1+\tfrac{\zeta_k}{\rho}\log\max\{1,h_0/H_{\mathrm w}\}$.
\item[(iii)] For $t\ge\max(t_{\mathrm w},t_{\mathrm{id}}+N)$ every
iteration contracts. Let $\rho^{\mathrm{act}}_t:=1-h_{t+1}/h_t$ (for
$h_t>0$) be the actual contraction and, for $j\ge1$,
\[
\underline\rho_j:=\min\Big\{1,\ 4\rho\max\Big[1,\frac{1+\beta_{\mathrm{flat}}(j-1)}{1+\eta(j-1)}\Big]\Big\},
\]
which is nondecreasing in $j$. Then $\rho^{\mathrm{act}}_t\ge\rho$, and
$\rho^{\mathrm{act}}_t\ge\underline\rho_{k'}$ when the iteration is case
(a) with an $\eta$-diverse, $\beta_{\mathrm{flat}}$-flat block of $k'$
pairs (under the two-sided schedule, the block of the first,
larger-gap side, which is the step the progress argument uses).
\end{enumerate}
\end{theorem}

\begin{proof}
Throughout, $s_i=\ip{g}{v_i}=-\ip{z}{v_i}$ on side $V$ and
$s_j=\ip{z}{w_j}$ on side $W$; we argue on side $V$, side $W$ being
symmetric with the maximum in place of the minimum.

\emph{Three facts.} Write $\varpi_t:=\norm{z_t-z^*}\le\sqrt{2h_t}$
by \eqref{eq:dist-h}. (F1) If $\varpi\le\tau/(6L)$ then for $i\notin F_V$ and
$j\in F_V$,
$s_j-s_i=\ip{z}{v_i-v_j}=\ip{z^*}{v_i-v_j}+\ip{z-z^*}{v_i-v_j}\ge\tau_i-\varpi L\ge5\tau/6$,
while for $i,j\in F_V$, $|s_i-s_j|\le\varpi L\le\tau/6$; in particular the
global argmax of $s$ is a face point. (F2) Let $\omega_t$ be the total
weight on non-face points of both sides. By convexity
$h_t\ge\ip{\nabla F(z^*)}{z_t-z^*}=\ip{z^*}{z_t-z^*}$; writing
$\vartheta_V:=\min_i\ip{z^*}{v_i}$ and $\vartheta_W:=\max_j\ip{z^*}{w_j}$, so that
$\ip{z^*}{v_i}=\vartheta_V+\tau_i$ and $\ip{z^*}{w_j}=\vartheta_W-\tau_j$, and using
$\norm{z^*}^2=\ip{z^*}{p^*-q^*}=\vartheta_V-\vartheta_W$ (the optimal weights live on the
faces, where the margins vanish),
$\ip{z^*}{z_t-z^*}=\sum_i\alpha_i\tau_i+\sum_j\beta_j\tau_j\ge\tau\,\omega_t$.
Hence $\omega_t\le h_t/\tau$. (F3) By \eqref{eq:dist-h} and (F2),
$h_t<H_{\mathrm{id}}$ implies $\varpi_t\le\tau/(6L)$, $\omega_t\le\tau/(72L^2)$,
$\omega_t\le1/12$ and $2h_t+R\sqrt{2h_t}\le\tau/6$.

\emph{(i) Identification.} Let $h_t<H_{\mathrm{id}}$ and suppose
$\operatorname{supp}\alpha\not\subset F_V$. The worst active donor $u$
minimises $s$ over the support, so by (F1) $u\notin F_V$. The away gap of
side $V$ is $\gA_V=\sum_i\alpha_i(s_i-s_u)\ge\sum_{i\in F_V}\alpha_i\cdot5\tau/6\ge(1-\omega_t)5\tau/6\ge55\tau/72$.
A side whose support lies in its face has
$\gPW\le\gA+\gFW\le\varpi L+(2h+R\sqrt{2h})\le\tau/3$ by (F1) and (F3), so the
side with the larger pairwise gap, which the schedule steps first, is a
side with a non-face support point whenever one exists. On that side
pair~1 is $(r^*,u)$ with $r^*\in F_V$ the global argmax, so
$\mathrm{gap}_1\ge5\tau/6$ and the uncapped transfer would be
$\mathrm{gap}_1/\norm{d_1}^2\ge5\tau/(6L^2)>\tau/(72L^2)\ge\omega_t\ge\alpha_u$:
pair~1 is capped and the iteration is case (b). The away step on $u$ has
derivative $\ip{g}{p-v_u}=\gA_V\ge55\tau/72$ and curvature
$\norm{p-v_u}^2\le L^2$, so its unconstrained length is at least
$55\tau/(72L^2)$, while its maximal length is
$\alpha_u/(1-\alpha_u)\le\tfrac{12}{11}\alpha_u\le\tau/(11L^2)$; the
away step is boundary-clipped, i.e.\ a drop of $u$, and it adds no index.
If instead neither side has a non-face support point, all donors are
face points and a non-face receiver $r$ has
$\mathrm{gap}=s_r-s_u\le-5\tau/6<0$ against any of them, so it is
discarded (the discard rule of Algorithm~\ref{alg:block}); the toward candidate is
$r^*\in F_V$. Under the two-sided schedule the second side steps only on
non-drop iterations, when neither side has a non-face point, and the same
argument applies. Thus for $t\ge t_{\mathrm{id}}$ (monotonicity of $h$)
the set of non-face support indices never grows and loses one member at
every iteration while it is nonempty, which proves (i); the bound on $N$
is the support count of Theorem~\ref{thm:linear}.

\emph{(ii) No late drops.} Let $t\ge\max(t_{\mathrm w},t_{\mathrm{id}}+N)$,
so $p_t\in\mathcal F_V$, $q_t\in\mathcal F_W$ and, by \eqref{eq:dist-h}
applied to the strict inequality $h_t<H_{\mathrm w}$,
$\varpi_t<\min\{\tau/(6L),\sigma\alpha_{\min}/2,\alpha_{\min}d_{\min}/2\}$.
By (A2), $\norm{\alpha_t-\alpha^*}_\infty\le \varpi_t/\sigma<\alpha_{\min}/2$,
so by (A1) every face index carries weight
$\alpha_{t,i}>\alpha^*_i-\alpha_{\min}/2\ge\alpha_{\min}/2>0$: the
supports equal the faces and every support weight exceeds
$\alpha_{\min}/2$. Any pair $(r,u)$ of face points has
$\mathrm{gap}_j=s_r-s_u=\ip{z-z^*}{v_u-v_r}\le\varpi_t\norm{d_j}$,
because $\ip{z^*}{\cdot}$ is constant on the face, and
$\norm{d_j}=\norm{v_r-v_u}\ge d_{\min}$ (the distance between two
vertices is at least the distance from one of them to the affine hull of
the others), so its uncapped transfer satisfies
\[
\gamma_j\ =\ \frac{\mathrm{gap}_j}{\norm{d_j}^2}\ \le\ \frac{\varpi_t}{\norm{d_j}}\ \le\ \frac{\varpi_t}{d_{\min}}\ <\ \frac{\alpha_{\min}}2\ <\ \alpha_{u_j}:
\]
no pair is capped, the iteration is case (a), no away step is taken,
and a block transfer removes strictly less than $\alpha_{u_j}$ from each
donor, so no index leaves the support. (If both faces are singletons
there are no pairs and nothing to show.) Drops at $t\ge t_{\mathrm w}$ are therefore among the
$N$ identification drops of (i), which occur at the consecutive iterations
$t_{\mathrm{id}},\dots,t_{\mathrm{id}}+N-1$; hence all drops occur before
$t_{\mathrm w}+N$, and the count of Theorem~\ref{thm:linear} over the
first $t_{\mathrm w}+N$ iterations gives the bound. The bound on
$t_{\mathrm w}$ is Theorem~\ref{thm:linear} with
$\log(1/(1-\rho))\ge\rho$.

\emph{(iii) Rate.} A non-drop iteration gains at least the amount of
Lemma~\ref{lem:good}, which with Lemma~\ref{lem:gsc} is $\rho h_t$. In
case (a) the guard takes $\max(\Delta_B,\Delta_1)$, which is at least
$\Delta_1$ (the guard) and, by Proposition~\ref{prop:ratio}, at least
$\Delta_1(1+\beta_{\mathrm{flat}}(k'-1))/(1+\eta(k'-1))$; hence at least
$\Delta_1\max[1,(1+\beta_{\mathrm{flat}}(k'-1))/(1+\eta(k'-1))]$. In case
(a) the proof of Lemma~\ref{lem:good} gives
$\Delta_1\ge(\gPW)^2/(8M^2)\ge2\delta^2h_t/(8M^2)=4\rho h_t$ by
Lemma~\ref{lem:gsc}. A gain cannot exceed $h_t$, whence the minimum
with~$1$. Monotonicity of $\underline\rho_j$ in $j$: the derivative of
$(1+\beta_{\mathrm{flat}}(j-1))/(1+\eta(j-1))$ in $j$ has the sign of
$\beta_{\mathrm{flat}}-\eta$, so the ratio is nondecreasing when
$\beta_{\mathrm{flat}}\ge\eta$, and when $\beta_{\mathrm{flat}}<\eta$ it
lies below $1$ and the outer maximum makes $\underline\rho_j$ constant.
\end{proof}

Theorem~\ref{thm:ident} says that drops are paid during identification
and never afterwards: the worst case $k$ drops per iteration of
Theorem~\ref{thm:linear} applies only up to $t_{\mathrm w}+N$, and the
$k$-dependence of the drop count enters only through the support that
is built up to then. This is exactly the shape of the measurements in
Figure~\ref{fig:diag}. The thresholds are those of active-set
identification and are pessimistic (Section~\ref{sec:ident-check}
measures by how much); the measured drop counts below are
$0.3$--$70\%$ of one class's final support.

\begin{corollary}[Conditional contraction and iteration count]\label{cor:work}
Run the algorithm with requested block size $k$. At iteration $t$ let
$k'_t\in\{1,\dots,k\}$ be the number of pairs retained by the block used
in the progress argument (the first side's block; $k'_t=1$ on case-(b)
and drop iterations, where only the leading pair is used). Let
$t_1(k):=\max(t_{\mathrm w},t_{\mathrm{id}}+N)$ as in
Theorem~\ref{thm:ident} and suppose every case-(a) iteration
$t\ge t_1(k)$ is $\eta$-diverse and $\beta_{\mathrm{flat}}$-flat. Then
for every $T\ge t_1(k)$,
\[
h_T\ \le\ h_{t_1(k)}\prod_{t=t_1(k)}^{T-1}\big(1-\rho^{\mathrm{act}}_t\big),\qquad
\rho^{\mathrm{act}}_t\ \ge\ \underline\rho_{k'_t}
\]
on case-(a) iterations (after $t_1(k)$ every iteration is case (a)
with no capped pair, by Theorem~\ref{thm:ident}(ii), so the convention
$k'_t=1$ matters only before $t_1(k)$). In particular, for a target $\varepsilon_h$ on the objective error
(distinct from the relative certificate tolerance $\varepsilon$ of
Section~\ref{sec:setting}; Corollary~\ref{cor:std} converts one into
the other): if $h_{t_1(k)}>\varepsilon_h$ and $k'_t\ge\underline k$ for
some $\underline k\in\{1,\dots,k\}$ on every case-(a) iteration after
$t_1(k)$, then $h_T\le\varepsilon_h$ for
$T=t_1(k)+\big\lceil\log(h_{t_1(k)}/\varepsilon_h)/\underline\rho_{\underline k}\big\rceil$
when $\underline\rho_{\underline k}<1$ (and $T=t_1(k)+1$ when
$\underline\rho_{\underline k}=1$); if $h_{t_1(k)}\le\varepsilon_h$ the
run is over at $t_1(k)$.
\end{corollary}
\begin{proof}
Theorem~\ref{thm:ident}(iii) with $k'$ read as $k'_t$ gives the product;
since $\underline\rho_j$ is nondecreasing in $j$,
$\underline\rho_{k'_t}\ge\underline\rho_{\underline k}$, and
$\log(1/(1-\underline\rho))\ge\underline\rho$ turns the product into the
iteration count.
\end{proof}

\paragraph{Timing model.} The corollary counts iterations; what they
cost is a matter of the fitted coefficients, not of a theorem. Let
$k^{\mathrm{ev}}_t\le2k$ be the total number of pairs the implementation
assembled and evaluated at iteration $t$ (both sides' blocks under the
two-sided schedule, including the pairs evaluated before an away drop)
and $N_{\mathrm{cols}}(k)$ the number of distinct columns computed over
the whole run. The cost model approximates the run's work by
\[
W(k)\ \approx\ c_{\mathrm{col}}N_{\mathrm{cols}}(k)+\sum_{t<T}\big(a+b\,k^{\mathrm{ev}}_t\big)
\ \lesssim\ c_{\mathrm{col}}N_{\mathrm{cols}}(k)+(a+2bk)\Big[t_1(k)+\Big\lceil\frac{\log(h_{t_1(k)}/\varepsilon_h)}{\underline\rho_{\underline k}}\Big\rceil\Big],
\]
the second expression using the iteration count of
Corollary~\ref{cor:work} and $k^{\mathrm{ev}}_t\le2k$. Here $b$ is a
cost per evaluated pair; the $b$ fitted in Section~\ref{sec:pay-meas}
is a cost per iteration and requested size $k$, which under the
two-sided schedule absorbs the factor two.
Corollary and model together describe one run. They do not by
themselves compare two block sizes: $N_{\mathrm{cols}}(k)$, $t_1(k)$ and
$h_{t_1(k)}$ all depend on the trajectory, and comparing the
multiplicative factors of two upper bounds is not a proof that one
algorithm is faster. What they do say is which factor aggregation can act
on. The per-iteration term after identification carries the factor
$(a+bk^{\mathrm{ev}}_t)/\underline\rho_{k'_t}$, and if
$k'_t\equiv k^{\mathrm{ev}}_t\equiv k$ (one-sided schedule, full blocks)
and $\beta_{\mathrm{flat}}=1$ this is proportional to
$(1+\eta(k-1))(a+bk)/k$, minimised at $k^*$ of
\eqref{eq:kstar}; the column term is not touched by the step rule at
all. Whether $N_{\mathrm{cols}}$, $t_1$ and $h_{t_1}$ are in fact
insensitive to $k$ is an empirical question, answered below for the
five cells (column counts within $20\%$ across $k$; drops and misses
concentrated in the same early deciles for every $k$). A comparative
speed-up theorem would need those three quantities bounded uniformly in
$k$, which we do not prove.

